\section{Introduction}

To a point $x$ in a complex Grassmannian $G(r,n)$ we associate two objects. The matroid $M_x$ whose bases are given by the non-zero Plücker coordinates of $x$ and the projective, toric subvariety of $G(r,n)$ given by the closure of the orbit of $x$ under the action of the algebraic torus $T= (\CC^*)^n$, denoted $\T$. The class of $\T$ in the Chow ring of $G(r,n)$ is an additive, valuative matroid invariant. Using results of Fink, Derksen, Speyer and Berget in \cite{DerksenFink}, \cite{FinkSpeyer} and \cite{BergetFink} we consider the unique valuative extension of this invariant to all matroids, not just those representable over $\CC$. We denote it by $\Sc(M) \in A^\bullet(G(r,n))$. The Chow ring $A^\bullet(G(r,n))$ is generated as a free abelian group by the Schubert cycles $\sigma_\lambda$ where $\lambda$ ranges over all partitions whose Young diagram is contained in the $r\times (n-r)$ rectangle. Hence $\Sc(M)$ can be expressed as the sum $\Sc(M) = \sum_\lambda d_\lambda(M) \sigma_\lambda$. We call the integers $d_\lambda(M)$ the Schubert coefficients of $M$.

The non-negativity of the Schubert coefficients was conjectured by Berget and Fink in \cite{BergetFink}. We explain how to compute the Schubert coefficients of a disconnected matroid in terms of the Schubert coefficients of its connected components in \cref{thm:disconnected}. And we compute the Schubert coefficients of sparse paving matroids and confirm their non-negativity, supporting the conjecture of \cite{BergetFink}.
\begin{thm}
    \label{thm:1}
    Let $M$ be a connected sparse paving matroid of rank $r$ on $[n]$ then $d_\lambda(M) = d_\lambda(U_{r,n})$ for all partitions $\lambda$ except for the complement of the hook partition $h^c$ for which $d_{h^c}(M) = \beta(M)$.
\end{thm}
\begin{corollary}
    \label{cor:1}
    Let $M$ be a connected sparse paving matroid. The Schubert coefficients $d_\lambda(M)$ are non-negative.
\end{corollary}
We suspect that sparse paving matroids are the only matroids such that $d_\lambda(M) = d_\lambda(U_{r,n})$ for all partitions $\lambda \neq h^c$, and make the following conjecture.
\begin{conj}
    \label{conj:sparsePaving}
    Let $M$ be a connected matroid of rank $r$ on $[n]$. $M$ is sparse paving if and only if $d_\lambda(M) = d_\lambda(U_{r,n})$ for all partitions $\lambda \neq h^c$.
\end{conj}

Section 2 consists of preliminaries on Grassmannians, matroids and valuations. Section 3 introduces the concept of Schubert coefficients, and recaps what is previously known. Section 4 is about Schubert coefficients of disconnected matroids. Section 5 contains the proof of \cref{thm:1} and some evidence for \cref{conj:sparsePaving}.